% Condensed in rewrite 2026-09-25: selected decoded examples (fig:segments), decoded-segment side analysis, Hamlyn border-restricted analysis and long protocol prose cut; tab:segments and a condensed protocol remain in App. D.7.
\subsection{Segmentation view of the forecasts}\label{app:segments}
\begin{figure}[h]
  \centering
  \figorpending{segments.pdf}{5.4cm}{Segmentation read out of the forecasts: SAM 2.1 masks, forecast masks, IoU vs horizon.}
  \caption[Segmentation view of the forecasts]{\ifdefined\segCaption\segCaption\else Segmentation view of the forecasts.\fi}
  \label{fig:segments}
\end{figure}
\IfFileExists{submission_folder/tables/generated/segments_decoded_text.tex}{\paragraph{Decoded forecasts.} To show the placement in pixels we decode every $k{=}10$ forecast with the trained feature decoder and segment the tool in the \emph{decoded} image with one rule for all models (top Grounding~DINO box for the dataset prompt $\rightarrow$ SAM~2.1; no segment if the detection score is below 0.3), then compare it with the SAM~2.1 mask of the true frame (patch-level IoU). On 668 moving DROID windows the decoded-segment IoU is 0.34 for \ours{}, 0.30 for Direct and 0.23 for AR (decoded true features: 0.54); paired gain of \ours{} +3.9 [+1.2, +6.6] IoU points vs.\ Direct and +10.3 [+7.4, +13.2] vs.\ AR (detection rate 83\%/83\%/80\%). On 509 moving Hamlyn windows the decoded-segment IoU is 0.39 for \ours{}, 0.38 for Direct and 0.38 for AR (decoded true features: 0.40); paired gain of \ours{} -1.0 [-3.8, +1.4] IoU points vs.\ Direct and -0.1 [-2.6, +2.3] vs.\ AR (detection rate 93\%/93\%/94\%). The examples in \cref{fig:segments} are selected by a fixed rule that favours \ours{}: moving windows where its decoded segment matches the target (IoU $\ge 0.6$ or placement error in its best quartile) while Direct and AR both miss it (IoU $\le 0.3$ or at least twice the placement error), ranked by the IoU margin over the better baseline, one per episode; they illustrate the effect, the averages above measure it. Every model is shown on the same window and the same frame.
}{}

\paragraph{Does the moving arm land in the right place? (Evaluation beyond feature error.)} Feature error does not show \emph{where} a forecast puts the
arm, so we read a segmentation out of every forecast (\cref{fig:segments}). SAM~2.1 \citep{ravi2025sam2} tracks the robot
arm (the instruments on Hamlyn) through the true future frames, and each forecast patch is labelled arm or background by
nearest-neighbour transfer from the observed frame alone (protocol below). \emph{Placement} is the distance between the
centroid of the cells a forecast labels as arm and the centroid of the true mask at $t{+}k$.

On the half of the
\segDroidWindows{} held-out DROID windows where the arm moves most, the mean placement error over $k$ is
\segDroidPlaceSW{}\,px for \ours{}, \segDroidPlaceDi{} for Direct, \segDroidPlaceAR{} for AR and \segDroidPlacePers{} for
persistence (\segDroidPlaceFrame{} frame): \ours{} removes \good{\segDroidPlaceRedAR{}\,px} relative to AR (95\% CI
\segDroidPlaceRedARCI{}) and \segDroidPlaceRedDi{}\,px relative to Direct (\segDroidPlaceRedDiCI{}). On Hamlyn
\ours{} places the instruments at \segHamlynPlaceSW{}\,px, Direct at \segHamlynPlaceDi{} and AR at \segHamlynPlaceAR{}\,px
(paired reductions \segHamlynPlaceRedDi{} \segHamlynPlaceRedDiCI{} and \segHamlynPlaceRedAR{} \segHamlynPlaceRedARCI{}\,px),
against \segHamlynPlacePers{}\,px for persistence.

Overlap
tells the same story: on moving DROID windows the arm read out of \ours{} has IoU \good{\segDroidMovGainDi{}} points above
Direct (\segDroidMovGainDiCI{}) and \good{\segDroidMovGainAR{}} above AR (\segDroidMovGainARCI{}); on Hamlyn the
gains are \segHamlynMovGainDi{} \segHamlynMovGainDiCI{} and \segHamlynMovGainAR{} \segHamlynMovGainARCI{} points
(\cref{tab:segments}). On Hamlyn, \segHamlynBorderExcluded{} reference masks do not touch the frame border and may
capture tissue rather than an instrument; restricted to the \segHamlynBorderWindows{} windows whose mask does
(\segHamlynBorderMoving{} moving), \ours{}, Direct and AR place the instruments at \segHamlynBorderPlaceSW{},
\segHamlynBorderPlaceDi{} and \segHamlynBorderPlaceAR{}\,px on moving windows (reductions \segHamlynBorderRedDi{}
\segHamlynBorderRedDiCI{} and \segHamlynBorderRedAR{} \segHamlynBorderRedARCI{}\,px), with IoU gains of
\segHamlynBorderGainDi{} \segHamlynBorderGainDiCI{} and \segHamlynBorderGainAR{} \segHamlynBorderGainARCI{} points. The examples in \cref{fig:segments}a are
decoded forecasts selected by the rule stated under \emph{Decoded forecasts} (\ours{} segment on target, Direct and AR
both off it); they illustrate the effect, and the claims rest on the averages.
\paragraph{Results.} \Cref{tab:segments} lists IoU (a) and placement error (b) per dataset and method. On DROID, \ours{} has the highest IoU at every reported horizon, on all and on
moving windows, with paired CIs that exclude zero against both Direct and AR. On Hamlyn, where the instruments move
less within ten steps, the three learned predictors are within half an IoU point of each other.
\begin{table}[h]
  \centering\scriptsize
  \caption{\textbf{Segmentation view of the forecasts.} (a) IoU of the foreground read out of each forecast
  (horizon-averaged, $k{=}5$ and $k{=}K$; all and moving windows). (b) Arm placement error (px of the displayed frame,
  DROID $320{\times}180$, Hamlyn $396{\times}224$; lower is better). Gains: moving windows, mean over $k$, paired 95\%
  CI. \good{Green}: in (a), paired 95\% CI of \ours{} minus \emph{both} Direct and AR excludes zero; in (b), CI excludes zero.}
  \label{tab:segments}
  \begin{tabular}{llccccccc}
    \toprule
    \multicolumn{9}{@{}l}{\textbf{(a) IoU} $\uparrow$} \\
    & & \multicolumn{3}{c}{all windows} & \multicolumn{2}{c}{moving windows} & \multicolumn{2}{c}{\ours{} gain, moving (IoU pts)} \\
    \cmidrule(lr){3-5}\cmidrule(lr){6-7}\cmidrule(lr){8-9}
    Dataset & Method & avg.\ over $k$ & $k{=}5$ & $k{=}K$ & avg.\ over $k$ & $k{=}K$ & vs Direct & vs AR \\
    \midrule
    \inputrows{submission_folder/tables/generated/segments_rows.tex}
    \midrule
    \multicolumn{9}{@{}l}{\textbf{(b) Placement error} (px) $\downarrow$} \\
    & & \multicolumn{2}{c}{all windows} & \multicolumn{3}{c}{moving windows} & \multicolumn{2}{c}{error removed by \ours{} (px)} \\
    \cmidrule(lr){3-4}\cmidrule(lr){5-7}\cmidrule(lr){8-9}
    Dataset & Method & avg.\ over $k$ & $k{=}K$ & avg.\ over $k$ & $k{=}5$ & $k{=}K$ & vs Direct & vs AR \\
    \midrule
    \inputrows{submission_folder/tables/generated/segments_place_rows.tex}
    \bottomrule
  \end{tabular}
\end{table}

\paragraph{Protocol.} We ask whether the moving part of the scene lands where it should in each forecast, without
decoding pixels. \emph{Reference masks.} On the observed frame $t$ (the last history frame) of every held-out window
(test split; the planning suites have no test split, so their validation split is used; stride 2 over all episodes, and an evenly spaced subset of 1{,}500 windows where there are more),
Grounding DINO \citep{liu2024groundingdino} detects a fixed per-dataset text prompt (DROID, Bridge, RT-1: ``a robot arm'';
Hamlyn: ``a surgical instrument'', up to two non-overlapping boxes; Language-Table: ``a robot arm. a block.'', up to two boxes; IWS: the manipulated T block, box or rope). Its
top box prompts SAM~2.1 (Hiera-L) \citep{ravi2025sam2}, which segments frame $t$ and \emph{tracks} the mask through the
true future frames $t{+}1,\dots,t{+}K$ with its video memory. Frames are segmented at their native aspect ratio; each
mask keeps its largest connected component and any component with $\ge$25\% of its area. For the rendered planning
suites a text prompt is unreliable, so an exact colour threshold of the agent/object is used instead.

\emph{Quality
control.} A window is kept only if the detection score is $\ge 0.3$, the mask covers 3--128 of the 256 patches at $t$,
the track is never lost and the mask area never changes by more than $3\times$ between consecutive frames; on DROID
this keeps \segDroidWindows{} of \segDroidConsidered{} windows (rejected: \segDroidRejNoDetection{} without a
confident detection, \segDroidRejFgAreaAtT{} with an arm too small or too large at $t$, \segDroidRejLostTrack{} lost
tracks, \segDroidRejAreaJump{} area jumps), and on Hamlyn \segHamlynWindows{} of \segHamlynConsidered{}.

\emph{Reading a mask out of a forecast.} Masks are area-averaged onto the $16{\times}16$ patch grid (a patch is
foreground if $\ge$50\% covered). Each patch of a forecast $\hat\Z_k$ is labelled by nearest-neighbour transfer
from the observed frame only: it is foreground if its most similar (cosine) foreground patch of $\Z_0$ is more similar
than its most similar background patch. This labeller was fixed in a DROID pilot as the one with the highest IoU when
applied to the \emph{true} future features, a criterion that involves no forecast (on DROID at $k{=}10$:
nearest-neighbour \segDroidOracleNn{}, $k$-means prototypes \segDroidOracleMulti{}, single prototypes
\segDroidOracleProto{}).

\emph{Metric.} IoU between the forecast mask and the tracked SAM mask of the true future frame,
per horizon and averaged over $k$, on all windows and on \emph{moving} windows (overlap of the masks at $t$ and $t{+}K$
below the median). Means pool all windows; their 95\% CIs resample episodes (all windows of each drawn episode); gains are means of paired per-episode differences with an episode bootstrap.
\emph{Placement.} The predicted centroid is the mean of the centres of the cells labelled arm; the true centroid is
the coverage-weighted mean of cell centres under the tracked SAM mask; distances are in pixels of the displayed
native-aspect frame (DROID \segDroidPlaceFrame{}, Hamlyn \segHamlynPlaceFrame{}). At a given $k$, windows in which any compared method labels no cell are dropped for all methods (paired; this
affects at most 1\% of DROID and 3\% of Hamlyn windows).

\emph{Examples} (\cref{fig:segments}) show each model's own decoded
$k{=}10$ forecast, selected by the fixed rule stated under \emph{Decoded forecasts} above: moving windows where the decoded
segment of \ours{} matches the target while Direct and AR both miss it, ranked by IoU margin, one per episode, and shown
only if the reference mask at $t$ touches the image border and the \ours{} segment overlaps the true mask (IoU
$\ge 0.3$). All methods are always shown on the same window and the same frame. The examples are not representative; averages and paired CIs are in
\cref{tab:segments}. Checkpoints are the final ones used in the main tables.
\paragraph{Caveats.} (i) The readout measures where a forecast puts content that \emph{looks like} the observed
foreground; it cannot credit content that was not visible at $t$. (ii) The labeller is imperfect: applied to the true
future features it reaches IoU \segDroidOracleNn{} on DROID at $k{=}10$, so absolute values are lower bounds on
localisation quality; on Hamlyn several forecasts slightly exceed this reference, because smoother forecasts can be
labelled more easily than the true features. (iii) SAM~2.1 occasionally includes a grasped object or tool-attached
structure in the tracked mask; QC removes lost or exploding tracks but
not such merges. On Hamlyn most rejected windows (\segHamlynRejFgAreaAtT{} of \segHamlynRejected{}) are close-ups in which
the instruments cover more than half of the frame. (iv) One seed per method.
